\documentclass[10pt,letterpaper]{article}
\usepackage[utf8]{inputenc}\usepackage[T1]{fontenc}\usepackage{lmodern}
\usepackage[spanish,es-noshorthands]{babel}\usepackage{csquotes}
\usepackage[margin=1.7cm]{geometry}\usepackage{amsmath,amssymb,xcolor,array,booktabs,enumitem,longtable}
\definecolor{guinda}{HTML}{6F1D46}
\setlength{\parindent}{0pt}\setlength{\parskip}{3pt}

\begin{document}
{\color{guinda}\large\textbf{Clave comentada — Examen del 2.º parcial, versión A}}\\
{\small Sesiones 4 a 8 (sin el laboratorio de la 8) e Influencers · 24 reactivos · 1 a 3 correctas · \textbf{documento del profesor}}
\vspace{4pt}\hrule\vspace{4pt}
\textbf{Calificación:} por reactivo $\max\{0,(A-E)/k\}$; cada reactivo vale $10/24$. Letras correctas en esta versión: a: 13, b: 13, c: 13, d: 13.
\subsection*{Resumen de claves}\begin{center}\footnotesize\begin{tabular}{@{}rllcl@{}}\toprule N.º & Bloque & Tema & Correctas & Dificultad\\\midrule
1 & Sesión 4 & Matriz de criterios: asistente de compras con IA (comercio agéntico) & \textbf{a, b, c} & media\\
2 & Sesión 4 & Predicciones frente a lo verificable (cookies, Gartner) & \textbf{a, c} & media\\
3 & Sesión 4 & Tendencia o estacionalidad en una serie (Figura 1) & \textbf{b, d} & alta\\
4 & Sesión 4 & Escucha social: nube de palabras y base legal & \textbf{a, b, c} & media\\
5 & Sesión 5 (Google Analytics) & Qué es y qué no es Google Analytics & \textbf{a, c} & media\\
6 & Sesión 5 (Google Analytics) & Identidad en Google Analytics: \texttt{client\_id} y User-ID & \textbf{a, b, c} & media\\
7 & Sesión 5 (Google Analytics) & Retención de datos y exportación en Google Analytics & \textbf{b, d} & media\\
8 & Sesión 5 (Google Analytics) & Lectura de un embudo por dispositivo (tabla) & \textbf{a, c, d} & alta\\
9 & Sesión 6 & Nudge frente a patrón oscuro & \textbf{b} & media\\
10 & Sesión 6 & Reglas de decisión: aditiva, conjuntiva y lexicográfica (tabla) & \textbf{a, b, d} & alta\\
11 & Sesión 6 & Racionalidad limitada, estrategias adaptativas y heurísticas & \textbf{a, d} & media\\
12 & Sesión 6 & Lectura cuidadosa de cifras del comercio digital en México & \textbf{b} & media\\
13 & Sesión 7 & Genealogía de AIDA y respaldo de los marcos & \textbf{b, c} & alta\\
14 & Sesión 7 & AIDAS: etapa, dato y métrica & \textbf{a, c, d} & media\\
15 & Sesión 7 & Comparación de escenarios con un modelo de recorrido (tabla) & \textbf{b, c} & alta\\
16 & Sesión 7 & Atribución: último clic frente a efecto de remoción & \textbf{a, d} & alta\\
17 & Sesión 8 (teoría) & Siete papeles en una misma compra (caso) & \textbf{c, d} & media\\
18 & Sesión 8 (teoría) & Elegir creador con datos: costo por interacción, por alcance y señales de alerta (tabla) & \textbf{a, b} & alta\\
19 & Sesión 8 (teoría) & Evidencia sobre influenciadores y la guía de PROFECO & \textbf{c, d} & media\\
20 & Sesión 8 (teoría) & Qué dato deja cada figura participante & \textbf{d} & media\\
21 & Influencers & Clasificar a un creador: tamaño, rol y plataforma & \textbf{c, d} & media\\
22 & Influencers & Benchmarks de engagement y modelos de pago & \textbf{a, c} & media\\
23 & Influencers & Comparar creadores: ER, CPE y CPM (tabla) & \textbf{b, d} & alta\\
24 & Influencers & Riesgos y transparencia en México & \textbf{a, b, d} & media\\
\bottomrule\end{tabular}\end{center}\subsection*{Justificación por reactivo}
\vspace{3pt}\noindent\textbf{1.} \textit{Matriz de criterios: asistente de compras con IA (comercio agéntico)} (reactivo base 1) — correctas: \textbf{a, b, c}\\[-1pt]
{\scriptsize Fuente: Sesión 4: «Frontera: el comercio agéntico» y «La intermediación puede sesgar la decisión»}
\begin{itemize}[leftmargin=1.2em,itemsep=0pt,topsep=1pt,parsep=0pt]
\item[\textbf{a)}] {\footnotesize \textbf{[correcta]} Modelo de negocio: la comisión alinea sus incentivos con cerrar la venta, no necesariamente con el mejor interés del consumidor.\\ \textit{Por qué:} Criterio 1: cada modelo (publicidad, suscripción, comisión) alinea incentivos de forma distinta; el asistente con comisión puede favorecer a quien paga más.}
\item[\textbf{b)}] {\footnotesize \textbf{[correcta]} Una capa con suficiente control sobre la interfaz puede pasar de facilitar la decisión a sesgarla, sin cambiar de tecnología.\\ \textit{Por qué:} Mathur et al. (2019), 11 000 sitios: la misma capa pasa de facilitar a sesgar.}
\item[\textbf{c)}] {\footnotesize \textbf{[correcta]} Intermediación de datos: la pregunta durable es quién controla la capa entre el consumidor y el punto de venta, y quién es dueño del dato de comportamiento.\\ \textit{Por qué:} Criterio 2 y diapositiva del comercio agéntico: no importa qué asistente tiene más usuarios, sino quién controla la capa de intermediación.}
\item[\textbf{d)}] {\footnotesize \textbf{[incorrecta]} Por ser un asistente de IA y no una plataforma, el criterio de modelo de negocio deja de aplicar y solo importa su exposición regulatoria.\\ \textit{Por qué:} Los cuatro criterios aplican a cualquier plataforma «presente o futura», incluida una capa de intermediación nueva.}
\end{itemize}
\vspace{3pt}\noindent\textbf{2.} \textit{Predicciones frente a lo verificable (cookies, Gartner)} (reactivo base 2) — correctas: \textbf{a, c}\\[-1pt]
{\scriptsize Fuente: Sesión 4: «Dos predicciones razonables…», «Predicciones frente a lo verificable», «Por qué el libro no da cifras de cuota»}
\begin{itemize}[leftmargin=1.2em,itemsep=0pt,topsep=1pt,parsep=0pt]
\item[\textbf{a)}] {\footnotesize \textbf{[correcta]} Las fuentes de industria se contradicen sobre si la búsqueda cayó; por eso lo no confirmable en fuente primaria no se reporta como dato duro.\\ \textit{Por qué:} Tabla de predicciones y «Criterio de verificación»: esa decisión es «la evidencia más honesta» del capítulo.}
\item[\textbf{b)}] {\footnotesize \textbf{[incorrecta]} Las predicciones fallaron porque Gartner es una fuente de menor jerarquía que la literatura arbitrada; el caso es una crítica a esa firma.\\ \textit{Por qué:} La diapositiva lo dice expresamente: «No es una crítica a Gartner»; ambas fuentes eran de jerarquía comparable y es una demostración metodológica.}
\item[\textbf{c)}] {\footnotesize \textbf{[correcta]} Las cookies de terceros siguen activas por defecto en Chrome, porque Google revirtió su plan.\\ \textit{Por qué:} Verificable a sept. 2026 según la tabla de la sesión.}
\item[\textbf{d)}] {\footnotesize \textbf{[incorrecta]} Como la caída de 25\,\% se anunció en un informe de industria de prestigio, debe reportarse como dato duro.\\ \textit{Por qué:} El libro no reporta como dato duro lo que no puede confirmar en fuente primaria.}
\end{itemize}
\vspace{3pt}\noindent\textbf{3.} \textit{Tendencia o estacionalidad en una serie (Figura 1)} (reactivo base 3) — correctas: \textbf{b, d}\\[-1pt]
{\scriptsize Fuente: Sesión 4: «Cómo se separa la tendencia de la temporada», «Interpretación de negocio», «Caso: la marca de electrodomésticos»}
\begin{itemize}[leftmargin=1.2em,itemsep=0pt,topsep=1pt,parsep=0pt]
\item[\textbf{a)}] {\footnotesize \textbf{[incorrecta]} La tendencia y la estacionalidad estimadas en las últimas semanas son tan confiables como las del centro de la serie.\\ \textit{Por qué:} En los extremos el promedio móvil es incompleto: la estimación es menos confiable (y la estacionalidad ahí se infla).}
\item[\textbf{b)}] {\footnotesize \textbf{[correcta]} Un promedio móvil de 52 semanas estima la tendencia porque, al abarcar un ciclo anual completo, la temporada se cancela.\\ \textit{Por qué:} Método del laboratorio; la estacionalidad es serie $-$ tendencia.}
\item[\textbf{c)}] {\footnotesize \textbf{[incorrecta]} El alza de unos 15 puntos en las semanas 101--104 es evidencia suficiente de que la campaña funciona.\\ \textit{Por qué:} Antes de atribuir un salto a la campaña: ¿pasó lo mismo en esta época el año pasado? Aquí sí (semanas 49--52, alza de ~14 puntos).}
\item[\textbf{d)}] {\footnotesize \textbf{[correcta]} El alza equivalente de las semanas 49--52, sin campaña, sugiere un patrón anual que podría explicar el salto.\\ \textit{Por qué:} Es la «firma de la estacionalidad»: el mismo salto un ciclo antes.}
\end{itemize}
\vspace{3pt}\noindent\textbf{4.} \textit{Escucha social: nube de palabras y base legal} (reactivo base 4) — correctas: \textbf{a, b, c}\\[-1pt]
{\scriptsize Fuente: Sesión 4: «Salida: nube de palabras», «Ejercicio… clasifiquen los 10 términos», «Ética: la escucha social no está exenta de la ley»}
\begin{itemize}[leftmargin=1.2em,itemsep=0pt,topsep=1pt,parsep=0pt]
\item[\textbf{a)}] {\footnotesize \textbf{[correcta]} La nube sirve para explorar; para interpretar hace falta clasificar los términos (positivo, fricción, neutro) con un trabajo adicional.\\ \textit{Por qué:} «Una nube de palabras explora; no interpreta.»}
\item[\textbf{b)}] {\footnotesize \textbf{[correcta]} Que los comentarios sean públicos no basta para justificar cualquier uso analítico posterior; hay que verificar la LFPDPPP vigente.\\ \textit{Por qué:} Diapositiva de ética de la sesión.}
\item[\textbf{c)}] {\footnotesize \textbf{[correcta]} Una palabra como \textit{precio} puede ser positiva (\enquote{precio justo}) o de fricción (\enquote{subió mucho}): la frecuencia sola no lo distingue.\\ \textit{Por qué:} La nube mezcla menciones positivas y de fricción que por tamaño no se distinguen.}
\item[\textbf{d)}] {\footnotesize \textbf{[incorrecta]} Como el análisis es agregado y no publica nombres, la LFPDPPP deja de aplicar al corpus.\\ \textit{Por qué:} El corpus contiene texto de personas identificables o potencialmente identificables; aplican las obligaciones de base legal de la Sesión 3.}
\end{itemize}
\vspace{3pt}\noindent\textbf{5.} \textit{Qué es y qué no es Google Analytics} (reactivo base 5) — correctas: \textbf{a, c}\\[-1pt]
{\scriptsize Fuente: GA4: «¿Qué es Google Analytics?» (definición de trabajo, en términos de la Unidad I, lo que NO es) y «Una advertencia de vigencia»}
\begin{itemize}[leftmargin=1.2em,itemsep=0pt,topsep=1pt,parsep=0pt]
\item[\textbf{a)}] {\footnotesize \textbf{[correcta]} Complementa, y no sustituye, a las encuestas, la escucha social y el CRM.\\ \textit{Por qué:} Lo que NO es: no cubre la señal declarada ni la conducta fuera de tu sitio.}
\item[\textbf{b)}] {\footnotesize \textbf{[incorrecta]} Como es una plataforma gratuita y consolidada, sus menús y nombres son estables y un tutorial de hace varios años sigue siendo una guía confiable.\\ \textit{Por qué:} La plataforma cambia (menús, nombres, límites); hay que revisar la fecha del tutorial y comparar con la ayuda vigente.}
\item[\textbf{c)}] {\footnotesize \textbf{[correcta]} Es una fuente de datos de primera parte y de señal conductual: registra lo que las personas hacen en un sitio web o una app propios.\\ \textit{Por qué:} Definición de trabajo: servicio gratuito que registra interacciones y las resume en informes; en términos de la Unidad I, datos de primera parte y señal conductual.}
\item[\textbf{d)}] {\footnotesize \textbf{[incorrecta]} Permite saber por qué las personas actúan como lo hacen, por lo que puede sustituir a las encuestas y a la escucha social.\\ \textit{Por qué:} No dice por qué actúan así (señal declarada) ni qué hacen fuera de tu sitio; complementa, no sustituye.}
\end{itemize}
\vspace{3pt}\noindent\textbf{6.} \textit{Identidad en Google Analytics: \texttt{client\_id} y User-ID} (reactivo base 6) — correctas: \textbf{a, b, c}\\[-1pt]
{\scriptsize Fuente: GA4: «¿Cómo reconoce GA4 a una persona?»; Sesión 3 (omnicanalidad, Verhoef et al., 2015)}
\begin{itemize}[leftmargin=1.2em,itemsep=0pt,topsep=1pt,parsep=0pt]
\item[\textbf{a)}] {\footnotesize \textbf{[correcta]} Si inicia sesión y la organización implementa User-ID, Google Analytics puede unir sus dispositivos en un solo usuario.\\ \textit{Por qué:} User-ID es un identificador propio para usuarios con inicio de sesión.}
\item[\textbf{b)}] {\footnotesize \textbf{[correcta]} El caso ilustra el problema de la identidad única que separa lo multicanal de lo omnicanal (Verhoef et al., 2015).\\ \textit{Por qué:} Lo dice la propia diapositiva: es el problema de identidad única de la Sesión 3.}
\item[\textbf{c)}] {\footnotesize \textbf{[correcta]} Aparece como tres usuarios, porque el \texttt{client\_id} (cookie \texttt{\_ga}) reconoce navegadores, no personas.\\ \textit{Por qué:} Tres dispositivos sin sesión iniciada = tres usuarios.}
\item[\textbf{d)}] {\footnotesize \textbf{[incorrecta]} La cookie \texttt{\_ga} caduca por defecto a las 24 horas, por lo que cada visita se cuenta como un usuario nuevo.\\ \textit{Por qué:} La caducidad por defecto es de 2 años (configurable; los navegadores pueden acortarla).}
\end{itemize}
\vspace{3pt}\noindent\textbf{7.} \textit{Retención de datos y exportación en Google Analytics} (reactivo base 7) — correctas: \textbf{b, d}\\[-1pt]
{\scriptsize Fuente: GA4: «¿Cuánto tiempo se conservan los datos?», «Aviso: la cuenta de demostración no deja exportar», «Precauciones con lo que sale de GA4»}
\begin{itemize}[leftmargin=1.2em,itemsep=0pt,topsep=1pt,parsep=0pt]
\item[\textbf{a)}] {\footnotesize \textbf{[incorrecta]} La retención de 14 meses borra los informes estándar anteriores a ese límite.\\ \textit{Por qué:} Solo afecta a las Exploraciones y los embudos; los informes estándar agregados no se ven afectados.}
\item[\textbf{b)}] {\footnotesize \textbf{[correcta]} La Exploración se corta en torno a los 14 meses, aunque el informe estándar agregado muestre todo el periodo.\\ \textit{Por qué:} Consecuencia explícita de la diapositiva de retención.}
\item[\textbf{c)}] {\footnotesize \textbf{[incorrecta]} El archivo exportado trae una fila por persona, por lo que sirve para identificar a cada cliente.\\ \textit{Por qué:} Los datos vienen agregados (no hay filas por persona); además está prohibido subir información personal identificable.}
\item[\textbf{d)}] {\footnotesize \textbf{[correcta]} En la cuenta de demostración no se puede exportar porque es de solo consulta; es una restricción de esa cuenta, no un error del usuario.\\ \textit{Por qué:} Aviso de la sesión.}
\end{itemize}
\vspace{3pt}\noindent\textbf{8.} \textit{Lectura de un embudo por dispositivo (tabla)} (reactivo base 8) — correctas: \textbf{a, c, d}\\[-1pt]
{\scriptsize Fuente: GA4: «Práctica 2: un embudo de compra» (fricción del Cap. 1) y «Exploraciones: siete técnicas»}
\begin{itemize}[leftmargin=1.2em,itemsep=0pt,topsep=1pt,parsep=0pt]
\item[\textbf{a)}] {\footnotesize \textbf{[correcta]} En ambos dispositivos el mayor abandono ocurre entre \texttt{view\_item} y \texttt{add\_to\_cart} (75\,\% en móvil; 65\,\% en escritorio).\\ \textit{Por qué:} Móvil: 3 000/12 000 = 25 \% avanza; escritorio: 2 800/8 000 = 35 \% avanza.}
\item[\textbf{b)}] {\footnotesize \textbf{[incorrecta]} Como el mayor abandono es el del pago, conviene rediseñar primero el \textit{checkout} en ambos dispositivos.\\ \textit{Por qué:} El mayor abandono es el primer paso, no el pago (\texttt{begin\_checkout}$\to$\texttt{purchase}: 40 \% y 30 \%).}
\item[\textbf{c)}] {\footnotesize \textbf{[correcta]} El paso \texttt{begin\_checkout}$\to$\texttt{purchase} avanza más en escritorio (70\,\%) que en móvil (60\,\%).\\ \textit{Por qué:} 900/1 500 = 60 \%; 1 372/1 960 = 70 \%.}
\item[\textbf{d)}] {\footnotesize \textbf{[correcta]} La conversión de punta a punta es 7.5\,\% en móvil y 17.15\,\% en escritorio.\\ \textit{Por qué:} 900/12 000 = 7.5 \%; 1 372/8 000 = 17.15 \%.}
\end{itemize}
\vspace{3pt}\noindent\textbf{9.} \textit{Nudge frente a patrón oscuro} (reactivo base 9) — correctas: \textbf{b}\\[-1pt]
{\scriptsize Fuente: Sesión 6: «Frontera: del nudge al patrón oscuro», «El criterio: información veraz y salida fácil»}
\begin{itemize}[leftmargin=1.2em,itemsep=0pt,topsep=1pt,parsep=0pt]
\item[\textbf{a)}] {\footnotesize \textbf{[incorrecta]} Un nudge se distingue de un patrón oscuro por el tamaño de su efecto en la conversión, no por la veracidad ni por la facilidad de salida.\\ \textit{Por qué:} El criterio verificable es información veraz y salida fácil, no la magnitud del efecto.}
\item[\textbf{b)}] {\footnotesize \textbf{[correcta]} Un aviso \enquote{quedan 2 piezas} que muestra siempre el mismo número, sin relación con el inventario, es un patrón oscuro.\\ \textit{Por qué:} Fila «Escasez» de la tabla: «Anuncia \enquote{quedan 2} con un número fijo».}
\item[\textbf{c)}] {\footnotesize \textbf{[incorrecta]} Según Luguri y Strahilevitz (2021), los patrones leves casi cuadruplicaron la probabilidad de conservar un servicio no pedido y los agresivos más que la duplicaron.\\ \textit{Por qué:} Está invertido: los leves más que duplicaron; los agresivos casi la cuadruplicaron.}
\item[\textbf{d)}] {\footnotesize \textbf{[incorrecta]} Un temporizador de oferta es siempre un patrón oscuro, aunque informe la fecha real en que termina la promoción.\\ \textit{Por qué:} Informar la fecha real de término es el nudge de la tabla; el patrón oscuro es el temporizador que se reinicia al recargar.}
\end{itemize}
\vspace{3pt}\noindent\textbf{10.} \textit{Reglas de decisión: aditiva, conjuntiva y lexicográfica (tabla)} (reactivo base 10) — correctas: \textbf{a, b, d}\\[-1pt]
{\scriptsize Fuente: Sesión 6: «Tres reglas, dos familias», laboratorio (utilidad aditiva $U_i=\sum w_j s_{ij}$, umbral, tolerancia)}
\begin{itemize}[leftmargin=1.2em,itemsep=0pt,topsep=1pt,parsep=0pt]
\item[\textbf{a)}] {\footnotesize \textbf{[correcta]} La regla aditiva elige A (7.3), aunque tenga el puntaje de calidad más bajo.\\ \textit{Por qué:} A=7.3; B=6.3; C=7.1; D=6.1.}
\item[\textbf{b)}] {\footnotesize \textbf{[correcta]} A gana con la aditiva y queda fuera con la conjuntiva porque la primera deja que un atributo compense a otro y la segunda exige un mínimo en cada uno.\\ \textit{Por qué:} Diferencia entre regla compensatoria y no compensatoria.}
\item[\textbf{c)}] {\footnotesize \textbf{[incorrecta]} La regla lexicográfica con prioridad en calidad elige D, porque su entrega es la mejor.\\ \textit{Por qué:} Con prioridad en calidad gana C (9); la entrega solo desempata si la calidad empata dentro de la tolerancia.}
\item[\textbf{d)}] {\footnotesize \textbf{[correcta]} La regla conjuntiva con umbral 5 deja pasar solo a B y a D.\\ \textit{Por qué:} A falla en calidad (4) y C falla en entrega (2); D pasa justo con 5.}
\end{itemize}
\vspace{3pt}\noindent\textbf{11.} \textit{Racionalidad limitada, estrategias adaptativas y heurísticas} (reactivo base 11) — correctas: \textbf{a, d}\\[-1pt]
{\scriptsize Fuente: Sesión 6: «2.1.3 Evaluación: casi nadie pondera todo»; Simon (1955), Payne et al. (1988), Tversky y Kahneman (1974)}
\begin{itemize}[leftmargin=1.2em,itemsep=0pt,topsep=1pt,parsep=0pt]
\item[\textbf{a)}] {\footnotesize \textbf{[correcta]} Descartará primero con una regla no compensatoria (p. ej., un mínimo por atributo) y comparará después las sobrevivientes con una compensatoria.\\ \textit{Por qué:} Estrategias adaptativas: ante muchas alternativas se descarta primero y se compara después (Payne et al., 1988).}
\item[\textbf{b)}] {\footnotesize \textbf{[incorrecta]} Ponderará todos los atributos de las 40 alternativas con una regla aditiva, porque la interfaz digital lo facilita.\\ \textit{Por qué:} Casi nadie pondera todo: con tiempo y atención limitados no se hace.}
\item[\textbf{c)}] {\footnotesize \textbf{[incorrecta]} Las heurísticas de representatividad, disponibilidad y anclaje no se ven afectadas por cómo la interfaz ordena y presenta las opciones.\\ \textit{Por qué:} La interfaz puede activarlas o mitigarlas.}
\item[\textbf{d)}] {\footnotesize \textbf{[correcta]} Buscará una opción suficientemente buena en lugar de la óptima (racionalidad limitada).\\ \textit{Por qué:} Simon (1955).}
\end{itemize}
\vspace{3pt}\noindent\textbf{12.} \textit{Lectura cuidadosa de cifras del comercio digital en México} (reactivo base 12) — correctas: \textbf{b}\\[-1pt]
{\scriptsize Fuente: Sesión 6: «Dimensión del comercio digital en México (2025)»}
\begin{itemize}[leftmargin=1.2em,itemsep=0pt,topsep=1pt,parsep=0pt]
\item[\textbf{a)}] {\footnotesize \textbf{[incorrecta]} El 17.7\,\% indica que casi una de cada cinco personas usuarias de internet compra en línea.\\ \textit{Por qué:} Es la participación en las ventas minoristas, no una proporción de personas.}
\item[\textbf{b)}] {\footnotesize \textbf{[correcta]} Que 37.3\,\% haya comprado en línea no implica que 62.7\,\% decida fuera de internet (puede haber \textit{webrooming}), y las dos fuentes miden cosas distintas, por lo que sus cifras no se suman.\\ \textit{Por qué:} «Lectura cuidadosa» de la diapositiva.}
\item[\textbf{c)}] {\footnotesize \textbf{[incorrecta]} El \textit{webrooming} es probar en tienda y comprar en línea; por eso la cifra de ENDUTIH subestima a quienes compran en tienda.\\ \textit{Por qué:} Esa es la definición de showrooming; webrooming = investigar en línea y comprar en tienda.}
\item[\textbf{d)}] {\footnotesize \textbf{[incorrecta]} Como 37.3\,\% compró en línea, 62.7\,\% decide su compra fuera de internet.\\ \textit{Por qué:} No: pudo investigar en línea y pagar en tienda (webrooming).}
\end{itemize}
\vspace{3pt}\noindent\textbf{13.} \textit{Genealogía de AIDA y respaldo de los marcos} (reactivo base 13) — correctas: \textbf{b, c}\\[-1pt]
{\scriptsize Fuente: Sesión 7: «El modelo más citado no tiene inventor claro», «Qué respaldo tiene cada marco»}
\begin{itemize}[leftmargin=1.2em,itemsep=0pt,topsep=1pt,parsep=0pt]
\item[\textbf{a)}] {\footnotesize \textbf{[incorrecta]} El acrónimo AIDA fue acuñado por Strong en 1925, al sistematizar la secuencia que Lewis formuló en 1898.\\ \textit{Por qué:} El acrónimo es de C. P. Russell (1921, \textit{Printers' Ink}); 1925 es el año desde el que circula la atribución a Lewis (Strong).}
\item[\textbf{b)}] {\footnotesize \textbf{[correcta]} El \textit{consumer decision journey} (Court et al., 2009) es investigación de una consultora y no fue arbitrado: es un marco útil, no un experimento neutral.\\ \textit{Por qué:} Cuadro «Qué respaldo tiene cada marco».}
\item[\textbf{c)}] {\footnotesize \textbf{[correcta]} La atribución de AIDA a Lewis (1898) circula desde 1925 y, según Iwamoto (2024), carece de evidencia primaria.\\ \textit{Por qué:} Se repitió un siglo sin consultar los documentos; Sheldon y Dukesmith lo teorizaron antes.}
\item[\textbf{d)}] {\footnotesize \textbf{[incorrecta]} La crítica al embudo lineal solo la sostienen consultoras y plataformas; no tiene respaldo académico independiente.\\ \textit{Por qué:} Sí lo tiene: Santos y Gonçalves (2021, revista arbitrada).}
\end{itemize}
\vspace{3pt}\noindent\textbf{14.} \textit{AIDAS: etapa, dato y métrica} (reactivo base 14) — correctas: \textbf{a, c, d}\\[-1pt]
{\scriptsize Fuente: Sesión 7: «Lectura digital de cada etapa» y «Tres advertencias por etapa»}
\begin{itemize}[leftmargin=1.2em,itemsep=0pt,topsep=1pt,parsep=0pt]
\item[\textbf{a)}] {\footnotesize \textbf{[correcta]} Acción: el \textit{checkout} deja la transacción; la métrica es la tasa de conversión, que por ser la mejor medida suele sobrevalorarse.\\ \textit{Por qué:} Advertencia de la etapa Acción.}
\item[\textbf{b)}] {\footnotesize \textbf{[incorrecta]} Atención: el clic y el tiempo en página son el dato; la métrica asociada es la tasa de clics.\\ \textit{Por qué:} Eso corresponde a Interés. En Atención el dato es la impresión y el tiempo de visualización; la métrica, alcance y frecuencia.}
\item[\textbf{c)}] {\footnotesize \textbf{[correcta]} Satisfacción: la calificación y la devolución son el dato; las métricas asociadas son la recompra y la recomendación.\\ \textit{Por qué:} Fila Satisfacción del cuadro.}
\item[\textbf{d)}] {\footnotesize \textbf{[correcta]} Deseo: guardar o agregar al carrito deja el evento «agregar al carrito»; la métrica asociada es el abandono de carrito.\\ \textit{Por qué:} Etapa más difícil de observar: los datos capturan solo su sombra.}
\end{itemize}
\vspace{3pt}\noindent\textbf{15.} \textit{Comparación de escenarios con un modelo de recorrido (tabla)} (reactivo base 15) — correctas: \textbf{b, c}\\[-1pt]
{\scriptsize Fuente: Sesión 7: laboratorio de Markov (casos A, B y C), «Interpretación: el costo de volver atrás», «Quien abre no cierra»; solo la lectura de resultados, sin ecuaciones}
\begin{itemize}[leftmargin=1.2em,itemsep=0pt,topsep=1pt,parsep=0pt]
\item[\textbf{a)}] {\footnotesize \textbf{[incorrecta]} Como el modelo supone que el siguiente paso depende de todo el historial de la persona, los resultados de los escenarios son exactos.\\ \textit{Por qué:} La cadena de Markov supone que el siguiente paso depende solo del estado actual; los resultados son tan buenos como sus supuestos.}
\item[\textbf{b)}] {\footnotesize \textbf{[correcta]} Comparar escenarios permite estimar el efecto de cambiar un supuesto (p. ej., más retrocesos) sin gastar presupuesto, pero depende de que las probabilidades de transición reflejen el comportamiento real.\\ \textit{Por qué:} Un modelo puede ser correcto en sus cálculos y engañoso en su uso.}
\item[\textbf{c)}] {\footnotesize \textbf{[correcta]} El escenario B convierte menos y tarda más que el A: volver atrás a evaluar baja la conversión y la retrasa; un embudo lineal que lo ignora sobreestimaría la conversión.\\ \textit{Por qué:} 30 \% $\to$ 22 \% y 4.0 $\to$ 5.5 contactos; el costo de volver atrás.}
\item[\textbf{d)}] {\footnotesize \textbf{[incorrecta]} En el escenario C el crédito de redes sociales sube con ambos modelos de atribución, por lo que el último clic y el efecto de remoción miden lo mismo.\\ \textit{Por qué:} Redes pasa de 4 \% a 6 \% con el último clic, pero de 20 \% a 31 \% con la remoción: miden cosas distintas y el último clic sigue subestimándola.}
\end{itemize}
\vspace{3pt}\noindent\textbf{16.} \textit{Atribución: último clic frente a efecto de remoción} (reactivo base 16) — correctas: \textbf{a, d}\\[-1pt]
{\scriptsize Fuente: Sesión 7: «Salida: ¿qué canal se lleva el crédito?», «Interpretación: quien abre no cierra», caso de la marca de ropa deportiva}
\begin{itemize}[leftmargin=1.2em,itemsep=0pt,topsep=1pt,parsep=0pt]
\item[\textbf{a)}] {\footnotesize \textbf{[correcta]} El efecto de remoción es una medida del modelo, no una prueba causal; para confirmarlo convendría un experimento (p. ej., apagar el canal en una región).\\ \textit{Por qué:} Pregunta guía 3 del caso: limitación del efecto de remoción como medida de causa.}
\item[\textbf{b)}] {\footnotesize \textbf{[incorrecta]} Si los cálculos del modelo de último clic son aritméticamente correctos, no puede llevar a una mala decisión.\\ \textit{Por qué:} Un modelo puede ser correcto en sus cálculos y engañoso en su uso.}
\item[\textbf{c)}] {\footnotesize \textbf{[incorrecta]} Trasladar el presupuesto a búsqueda es la decisión más segura, porque la acción es la etapa mejor medida y la más cercana a la causa de la compra.\\ \textit{Por qué:} La acción es la mejor medida y por eso la más sobrevalorada; presupuestar con el último clic recorta el canal que alimenta a los demás.}
\item[\textbf{d)}] {\footnotesize \textbf{[correcta]} Redes sociales suele abrir recorridos que otros canales cierran; el último clic no ve ese papel.\\ \textit{Por qué:} «Quien abre no cierra»: redes casi nunca es el último contacto.}
\end{itemize}
\vspace{3pt}\noindent\textbf{17.} \textit{Siete papeles en una misma compra (caso)} (reactivo base 17) — correctas: \textbf{c, d}\\[-1pt]
{\scriptsize Fuente: Sesión 8: «Valeria y la consola», «Siete papeles, no siete personas», «2.3.1…2.3.5», «Qué dato deja cada figura»}
\begin{itemize}[leftmargin=1.2em,itemsep=0pt,topsep=1pt,parsep=0pt]
\item[\textbf{a)}] {\footnotesize \textbf{[incorrecta]} El \textit{streamer} es el decisor, porque su reseña fue lo que más pesó en la compra.\\ \textit{Por qué:} El streamer es el influenciador (quien orienta); el decisor es quien determina si se compra, qué y cuándo.}
\item[\textbf{b)}] {\footnotesize \textbf{[incorrecta]} El abuelo es el comprador, y sus datos de transacción describen al usuario del producto.\\ \textit{Por qué:} Los datos de compra describen al comprador; los de uso, al usuario.}
\item[\textbf{c)}] {\footnotesize \textbf{[correcta]} Diego es usuario y consumidor, y al publicar el \textit{unboxing} también es prosumidor.\\ \textit{Por qué:} Prosumidor (Toffler, 1980): también produce contenido, como reseñas o \textit{unboxing}.}
\item[\textbf{d)}] {\footnotesize \textbf{[correcta]} La madre es la decisora y, al comparar reseñas, también ejerce el papel de evaluadora.\\ \textit{Por qué:} Una persona puede reunir varios papeles; decisor = si, qué y cuándo; evaluador = compara y pondera.}
\end{itemize}
\vspace{3pt}\noindent\textbf{18.} \textit{Elegir creador con datos: costo por interacción, por alcance y señales de alerta (tabla)} (reactivo base 18) — correctas: \textbf{a, b}\\[-1pt]
{\scriptsize Fuente: Influencers: «Tres métricas para comparar creadores», «Ejemplo numérico», «Cuatro riesgos» (audiencia inflada: engagement muy superior a 10 %)}
\begin{itemize}[leftmargin=1.2em,itemsep=0pt,topsep=1pt,parsep=0pt]
\item[\textbf{a)}] {\footnotesize \textbf{[correcta]} Z tiene el \textit{engagement rate} más alto (11\,\%) y el menor costo por interacción (\$1.82), pero un engagement tan alto justifica auditar su audiencia antes de elegirlo.\\ \textit{Por qué:} Z: 9 900/90 000 = 11 \%; 18 000/9 900 = \$1.82. Un engagement muy superior a 10 \% en Instagram es señal de alerta de audiencia inflada.}
\item[\textbf{b)}] {\footnotesize \textbf{[correcta]} Si el objetivo es notoriedad, Y es más eficiente que X por su menor costo por mil seguidores; si el objetivo es interacción, X es más eficiente que Y por su menor costo por interacción.\\ \textit{Por qué:} Y: \$120 vs. X: \$200 por mil seguidores; X: \$5.00 vs. Y: \$12.00 por interacción. Primero el objetivo, luego el creador.}
\item[\textbf{c)}] {\footnotesize \textbf{[incorrecta]} Y tiene un \textit{engagement rate} mayor que X porque suma más interacciones totales (3\,000 frente a 1\,200).\\ \textit{Por qué:} ER de X = 1 200/30 000 = 4 \%; ER de Y = 3 000/300 000 = 1 \%: el ER se calcula sobre los seguidores, no sobre el total de interacciones.}
\item[\textbf{d)}] {\footnotesize \textbf{[incorrecta]} Un \textit{engagement rate} mayor siempre indica una audiencia de mejor calidad, por lo que Z debe elegirse.\\ \textit{Por qué:} Un engagement excesivamente alto puede indicar bots o \textit{pods}; no es sinónimo de calidad.}
\end{itemize}
\vspace{3pt}\noindent\textbf{19.} \textit{Evidencia sobre influenciadores y la guía de PROFECO} (reactivo base 19) — correctas: \textbf{c, d}\\[-1pt]
{\scriptsize Fuente: Sesión 8: «2.3.2 Influenciador: más que seguidores», «Frontera: comercio social, live commerce y creadores»}
\begin{itemize}[leftmargin=1.2em,itemsep=0pt,topsep=1pt,parsep=0pt]
\item[\textbf{a)}] {\footnotesize \textbf{[incorrecta]} La Guía de Publicidad para Influencers de PROFECO (2023) es una ley específica con régimen sancionador propio.\\ \textit{Por qué:} Es una guía de recomendaciones con fundamento en la LFPC, no una ley específica.}
\item[\textbf{b)}] {\footnotesize \textbf{[incorrecta]} Las cifras de interacción por tamaño de audiencia (${\sim}$2.7\,\% nano; $<$1\,\% macro) provienen de un informe primario ya verificado.\\ \textit{Por qué:} Son fuentes secundarias de industria; no se localizó el informe primario (orden de magnitud).}
\item[\textbf{c)}] {\footnotesize \textbf{[correcta]} Los metaanálisis de Leung et al. (2022) y de Barari et al. (2025) convergen en que el efecto opera por credibilidad y atractivo percibidos.\\ \textit{Por qué:} 1 531 tamaños de efecto de 251 estudios y 571 de 71 artículos.}
\item[\textbf{d)}] {\footnotesize \textbf{[correcta]} Un influenciador pequeño puede ser un buen punto de contacto de confianza, pero es una hipótesis que cada marca debe contrastar con sus datos.\\ \textit{Por qué:} Cierre de la diapositiva de comercio social.}
\end{itemize}
\vspace{3pt}\noindent\textbf{20.} \textit{Qué dato deja cada figura participante} (reactivo base 20) — correctas: \textbf{d}\\[-1pt]
{\scriptsize Fuente: Sesión 8: «Qué dato deja cada figura», «2.3.3--2.3.5», «¿Dónde buscar a quien orienta la decisión?»}
\begin{itemize}[leftmargin=1.2em,itemsep=0pt,topsep=1pt,parsep=0pt]
\item[\textbf{a)}] {\footnotesize \textbf{[incorrecta]} Si el evaluador es un asistente de IA, deja de importar con qué regla pondera, porque la configura el usuario.\\ \textit{Por qué:} Vuelve la pregunta: ¿con qué regla de decisión pondera y quién la configuró?}
\item[\textbf{b)}] {\footnotesize \textbf{[incorrecta]} El 97.6\,\% de uso de internet de los jóvenes de 15 a 24 años (ENDUTIH 2025) indica a quién debe dirigirse el anuncio.\\ \textit{Por qué:} El dato dice dónde buscar, no a quién dirigirse.}
\item[\textbf{c)}] {\footnotesize \textbf{[incorrecta]} El decisor es la figura mejor medida, porque su elección siempre queda registrada como transacción.\\ \textit{Por qué:} De la transacción se desprende el comprador; del decisor, a menudo ningún dato visible.}
\item[\textbf{d)}] {\footnotesize \textbf{[correcta]} Los datos de uso describen al usuario y los de compra al comprador; del decisor, a menudo no hay ningún dato visible.\\ \textit{Por qué:} Tabla «Qué dato deja cada figura».}
\end{itemize}
\vspace{3pt}\noindent\textbf{21.} \textit{Clasificar a un creador: tamaño, rol y plataforma} (reactivo base 21) — correctas: \textbf{c, d}\\[-1pt]
{\scriptsize Fuente: Influencers: «¿Qué hace que alguien sea influencer?», «La pirámide», «El tamaño no lo dice todo: seis perfiles»}
\begin{itemize}[leftmargin=1.2em,itemsep=0pt,topsep=1pt,parsep=0pt]
\item[\textbf{a)}] {\footnotesize \textbf{[incorrecta]} Los cortes por seguidores (nano $<$10 mil; micro 10--100 mil; macro 100 mil--1 M; mega $>$1 M) son universales.\\ \textit{Por qué:} Son cortes de Influencer Marketing Hub; otras fuentes ubican el macro desde 250 mil.}
\item[\textbf{b)}] {\footnotesize \textbf{[incorrecta]} La relación parasocial es de ida y vuelta: el creador conoce a cada seguidor, y por eso la recomendación se siente de confianza.\\ \textit{Por qué:} Es cercanía de un solo sentido, como con un amigo que no sabe que existes.}
\item[\textbf{c)}] {\footnotesize \textbf{[correcta]} Una médica con 8\,000 seguidores que explica temas de salud es nano por tamaño y experta (líder de opinión) por rol.\\ \textit{Por qué:} Nano: menos de 10 mil seguidores; experto: pocos seguidores, mucha credibilidad.}
\item[\textbf{d)}] {\footnotesize \textbf{[correcta]} Un mismo creador puede ser micro por tamaño, experto por rol y afiliado por modelo de pago.\\ \textit{Por qué:} Diapositiva «seis perfiles».}
\end{itemize}
\vspace{3pt}\noindent\textbf{22.} \textit{Benchmarks de engagement y modelos de pago} (reactivo base 22) — correctas: \textbf{a, c}\\[-1pt]
{\scriptsize Fuente: Influencers: «El engagement baja al crecer… pero cuidado con el dato», «¿Cómo se les paga?»}
\begin{itemize}[leftmargin=1.2em,itemsep=0pt,topsep=1pt,parsep=0pt]
\item[\textbf{a)}] {\footnotesize \textbf{[correcta]} Casi todos los \textit{benchmarks} vienen de empresas del sector y de muestras distintas, por lo que se compara dentro de la misma plataforma, nicho y mercado.\\ \textit{Por qué:} «Lección de datos».}
\item[\textbf{b)}] {\footnotesize \textbf{[incorrecta]} En el modelo de comisión o afiliado se paga una tarifa fija por publicación, independiente de las ventas.\\ \textit{Por qué:} Comisión/afiliado = porcentaje por venta con código o enlace; la tarifa fija es otro modelo.}
\item[\textbf{c)}] {\footnotesize \textbf{[correcta]} En TikTok, una muestra de 1\,500 creadores mostró que el engagement no cae de forma constante con el tamaño (nano 7.0\,\%, micro 3.1\,\%, macro 3.9\,\%, mega 6.5\,\%).\\ \textit{Por qué:} Dato reportado en la sesión.}
\item[\textbf{d)}] {\footnotesize \textbf{[incorrecta]} Las tarifas por publicación en México (p. ej., \$300--\$2\,000 para nano) forman un tarifario oficial.\\ \textit{Por qué:} Son estimaciones de un blog (2026): referencias, no tarifario oficial.}
\end{itemize}
\vspace{3pt}\noindent\textbf{23.} \textit{Comparar creadores: ER, CPE y CPM (tabla)} (reactivo base 23) — correctas: \textbf{b, d}\\[-1pt]
{\scriptsize Fuente: Influencers: «Tres métricas para comparar creadores» y «Ejemplo numérico: ¿cuál conviene?»}
\begin{itemize}[leftmargin=1.2em,itemsep=0pt,topsep=1pt,parsep=0pt]
\item[\textbf{a)}] {\footnotesize \textbf{[incorrecta]} Si el objetivo es conversión, B es la mejor opción por su menor CPM.\\ \textit{Por qué:} El CPM mide costo por alcance potencial; B tiene el CPE más alto (\$12). Primero el objetivo, luego el creador.}
\item[\textbf{b)}] {\footnotesize \textbf{[correcta]} A y C tienen el mismo costo por interacción (\$5.00), aunque A tiene mayor \textit{engagement rate}.\\ \textit{Por qué:} A: 8 000/1 600 = 5; C: 18 000/3 600 = 5.}
\item[\textbf{c)}] {\footnotesize \textbf{[incorrecta]} C tiene el \textit{engagement rate} más alto (3\,\%).\\ \textit{Por qué:} A: 1 600/40 000 = 4 \%; C: 3 600/120 000 = 3 \%; B: 1 \%.}
\item[\textbf{d)}] {\footnotesize \textbf{[correcta]} B es el más barato por cada mil seguidores (\$120), pero tiene el \textit{engagement rate} más bajo de los tres (1\,\%).\\ \textit{Por qué:} A: \$200; B: 60 000/500 000×1 000 = \$120; C: \$150.}
\end{itemize}
\vspace{3pt}\noindent\textbf{24.} \textit{Riesgos y transparencia en México} (reactivo base 24) — correctas: \textbf{a, b, d}\\[-1pt]
{\scriptsize Fuente: Influencers: «Cuatro riesgos que debes saber ver», «Transparencia en México», «El mercado mexicano: lo que sí y lo que no sabemos»}
\begin{itemize}[leftmargin=1.2em,itemsep=0pt,topsep=1pt,parsep=0pt]
\item[\textbf{a)}] {\footnotesize \textbf{[correcta]} Según la guía de PROFECO (2023), el contenido pagado debe identificarse como publicidad; es una guía de recomendaciones y su fuerza vinculante sigue en discusión.\\ \textit{Por qué:} Diapositiva «Transparencia en México» y sus matices.}
\item[\textbf{b)}] {\footnotesize \textbf{[correcta]} La publicidad no declarada puede ser una infracción y, además, destruye la confianza que se pagó por tener.\\ \textit{Por qué:} Riesgo 4.}
\item[\textbf{c)}] {\footnotesize \textbf{[incorrecta]} La cifra de IAB México de 3.3 millones de creadores equivale a la inversión publicitaria oficial en influencers.\\ \textit{Por qué:} Son creadores de contenido, no inversión; no existe todavía una cifra oficial de inversión en influencers en México.}
\item[\textbf{d)}] {\footnotesize \textbf{[correcta]} Un engagement muy superior a 10\,\% en Instagram y comentarios genéricos repetidos son señales de audiencia inflada (bots, \textit{pods}).\\ \textit{Por qué:} Riesgo 1: audiencia inflada.}
\end{itemize}
\end{document}